% Optional additional section. Uses G_{n,k} from lerch_boundary_notes.tex.

\section{The logarithmic singularity at the Stieltjes endpoint}
\label{lerchboundary:sec:abel}

The opposite endpoint $\rho=1$ is continuous for the differentiated
Lerch family, but it is not infinitely differentiable in its deformation
parameter.  The following elementary Abel estimate identifies the first
singular term and transfers it to every simple endpoint zero.  Its
logarithmic structure is the classical one for power-logarithm
Dirichlet coefficients; the point here is its precise normalization
and its consequence for zero branches.

Put $\tau=-\log\rho$ and $L=\log(1/\tau)$, so that
$\rho\uparrow1$ corresponds to $\tau\downarrow0$.  For
$0\leq j\leq k-1$ define the absolutely convergent moments
\begin{equation}\label{lerchboundary:eq:moments}
 \mu_j(a)=\sum_{m=0}^{\infty}m^j f_{n,k}(a+m),
 \qquad
 H(a,\tau)=\sum_{j=0}^{k-1}\frac{(-\tau)^j}{j!}\mu_j(a).
\end{equation}
For $j=0$ use $m^0=1$, including at $m=0$.

\begin{theorem}[The first singular Abel coefficient]
\label{lerchboundary:thm:abel}
Fix $n\geq1$, $k\geq1$, and a compact interval $I\subset(0,\infty)$.
Uniformly for $a\in I$, as $\tau\downarrow0$,
\begin{equation}\label{lerchboundary:eq:abel}
 G_{n,k}(a,e^{-\tau})=H(a,\tau)
 +\frac{(-1)^{n+k}}{k!(n+1)}\tau^k L^{n+1}
 +O_{n,k,I}(\tau^k L^n).
\end{equation}
The family has continuous one-sided derivatives in $\tau$ through
order $k-1$ at zero. Its next derivative satisfies
\begin{equation}\label{lerchboundary:eq:abel-derivative}
 \partial_\tau^kG_{n,k}(a,e^{-\tau})
 =\frac{(-1)^{n+k}}{n+1}L^{n+1}+O_{n,k,I}(L^n).
\end{equation}
In particular it has no finite continuous $k$th derivative there.
The leading singular coefficient is independent of $a$.
\end{theorem}

\begin{proof}
Monicity of $P_{n,k}$ and a uniform expansion at large $m$ give
\begin{equation}\label{lerchboundary:eq:tail-coefficient}
 f_{n,k}(a+m)
 =\frac{(-1)^n(\log m)^n}{m^{k+1}}
 +O_{n,k,I}\left(\frac{(1+\log m)^{n-1}}{m^{k+1}}\right),
 \qquad m\geq1.
\end{equation}
Each $a$-derivative adds a factor $O(m^{-1})$, apart from bounded-degree
logarithmic factors.  All moments in \eqref{lerchboundary:eq:moments}
therefore converge normally on positive compact $a$-sets, including
any fixed number of $a$-derivatives. Differentiation in $\tau$ up to
order $k-1$ is justified by dominated convergence.

Write
$R_{k-1}(x)=e^{-x}-\sum_{j=0}^{k-1}(-x)^j/j!$. For $0\leq x\leq1$,
\[
 R_{k-1}(x)=\frac{(-1)^k x^k}{k!}+O_k(x^{k+1}),
\]
and for $x\geq1$, $|R_{k-1}(x)|\leq C_k x^{k-1}$.
Subtract $H$ from the defining series for $G$ and split the sum at
$m=1/\tau$. The leading part of the lower range is
\[
 \frac{(-1)^{n+k}\tau^k}{k!}
 \sum_{1\leq m\leq1/\tau}\frac{(\log m)^n}{m}
 =\frac{(-1)^{n+k}\tau^k}{k!(n+1)}L^{n+1}+O_{n,k}(\tau^k).
\]
The Taylor remainder in that range is
$O(\tau^{k+1}\sum_{m\leq1/\tau}(1+\log m)^n)
=O(\tau^k L^n)$.  The lower-order coefficient term in
\eqref{lerchboundary:eq:tail-coefficient} contributes
$O(\tau^k\sum_{m\leq1/\tau}(1+\log m)^{n-1}/m)
=O(\tau^kL^n)$. Finally the upper range is bounded by
\[
 C\tau^{k-1}\sum_{m>1/\tau}\frac{(1+\log m)^n}{m^2}
 =O(\tau^kL^n).
\]
These bounds prove \eqref{lerchboundary:eq:abel}, uniformly in $a$.

To avoid differentiating a remainder, compute the next derivative
directly for $\tau>0$:
\[
 \partial_\tau^kG_{n,k}(a,e^{-\tau})
 =(-1)^k\sum_{m\geq1}m^k e^{-m\tau}f_{n,k}(a+m).
\]
Equation~\eqref{lerchboundary:eq:tail-coefficient}, the same cutoff,
and $|e^{-x}-1|\leq x$ for $0\leq x\leq1$ give
\[
 \sum_{m\geq1}\frac{e^{-m\tau}(\log m)^n}{m}
 =\frac{L^{n+1}}{n+1}+O_n(L^n).
\]
The lower logarithmic powers contribute $O(L^n)$.
This proves \eqref{lerchboundary:eq:abel-derivative} independently
of differentiation of \eqref{lerchboundary:eq:abel}.
\end{proof}

\begin{theorem}[Endpoint regularity and singular displacement of simple zeros]
\label{lerchboundary:thm:abel-root}
Let $r>0$ be a simple zero of $G_{n,k}(\cdot,1)$, and write
$q=\partial_aG_{n,k}(r,1)\neq0$.
There is a unique nearby real zero $a(\tau)$ of
$G_{n,k}(\cdot,e^{-\tau})$ for all sufficiently small $\tau\geq0$.
Let $\alpha(\tau)$ be the real analytic zero of $H(\cdot,\tau)$
through $\alpha(0)=r$. Then
\begin{equation}\label{lerchboundary:eq:abel-root}
 a(\tau)=\alpha(\tau)
 -\frac{(-1)^{n+k}}{k!(n+1)q}\tau^k L^{n+1}
 +O_{n,k,r}(\tau^kL^n).
\end{equation}
The branch is $C^{k-1}$, but not $C^k$, at the endpoint, in either
$\tau$ or $\rho$. Moreover
\begin{equation}\label{lerchboundary:eq:abel-root-derivative}
 \frac{d^k a}{d\tau^k}
 =-\frac{(-1)^{n+k}}{(n+1)q}L^{n+1}+O_{n,k,r}(L^n).
\end{equation}
\end{theorem}

\begin{proof}
Uniform convergence of $G$ and $G_a$, together with $q\neq0$, gives
one and only one zero near $r$, continuously in $\tau\geq0$.
The ordinary analytic implicit function theorem applies to the finite
polynomial $H$, since $H_a(r,0)=q$.
Evaluating \eqref{lerchboundary:eq:abel} at $\alpha(\tau)$ and using a
uniform lower bound for $|G_a|$ first gives
$a-\alpha=O(\tau^kL^{n+1})$.  Taylor expansion in $a$, followed by
$H_a(\alpha(\tau),\tau)=q+O(\tau)$, gives
\eqref{lerchboundary:eq:abel-root}; the quadratic error in the
displacement and the $O(\tau)$ change of the denominator are both
$O(\tau^kL^n)$.

The convergent differentiated series show that $G$ has all mixed
derivatives of total order at most $k-1$ continuously at the endpoint.
Implicit differentiation, or the one-sided implicit function theorem,
therefore gives $a\in C^{k-1}$. Differentiate the implicit identity
$k$ times for $\tau>0$.  Every mixed derivative involving at least
one $a$-derivative remains bounded, because that derivative improves
the spectral tail by one inverse power of $m$. All derivatives of
$a$ of order below $k$ have finite limits. Thus the only unbounded
term other than $G_a a^{(k)}$ is the pure $\tau$ derivative in
\eqref{lerchboundary:eq:abel-derivative}. Solving for $a^{(k)}$ proves
\eqref{lerchboundary:eq:abel-root-derivative}. In this step replacing
$G_a(a(\tau),e^{-\tau})$ by $q$ incurs only $O(L^n)$:
its difference is $O(\tau+|a(\tau)-r|)$, because $G_{a\tau}$ is
bounded at the endpoint, and
$|a(\tau)-r|=O(\tau+\tau^kL^{n+1})$.
The leading coefficient is nonzero, so the $k$th derivative has no
finite endpoint limit.  Finally $\tau=-\log\rho$ is a smooth local
change of variable with nonzero derivative at $\rho=1$.
\end{proof}

For the two branches of Theorem~\ref{lerchboundary:thm:n2global}, this
gives a particularly visible endpoint effect. If $r_i=a_i(1)$ and
$q_i=\partial_aG_{2,1}(r_i,1)$, then $q_1<0<q_2$ and
\[
 a_i(\rho)=r_i+\frac{\tau}{3q_i}L^3+O(\tau L^2),\qquad
 \frac{da_i}{d\rho}=-\frac{1}{3q_i}L^3+O(L^2).
\]
Thus $a_1'(\rho)\to+\infty$ and $a_2'(\rho)\to-\infty$ as
$\rho\uparrow1$.  The nonmonotonicity of $a_1$ already follows from
the elementary endpoint and the rational Stieltjes signs; this Abel
calculation supplies its quantitative endpoint behavior.
